\documentclass[11pt,a4paper]{article}
\usepackage[english,greek]{babel}
\usepackage[utf8]{inputenc}
\usepackage{nimbusserif}
\usepackage[T1]{fontenc}
\usepackage[left=1.50cm, right=1.50cm, top=2.00cm, bottom=2.00cm]{geometry}
\usepackage{amsmath}
\let\myBbbk\Bbbk
\let\Bbbk\relax
\usepackage[amsbb,subscriptcorrection,zswash,mtpcal,mtphrb,mtpfrak]{mtpro2}
\usepackage{graphicx,multicol,multirow,enumitem,tabularx,mathimatika,gensymb,venndiagram,hhline,longtable,tkz-euclide,fontawesome5,eurosym,tcolorbox,tabularray}
\usepackage[explicit]{titlesec}
\tcbuselibrary{skins,theorems,breakable}
\usetikzlibrary{fillbetween}

% Tasks Package for Horizontal MCQs
\usepackage{tasks}
\settasks{label-format=\bfseries, label-offset=1em}

\pgfplotsset{compat=1.18}
\pgfkeys{/pgfplots/aks_on/.style={axis lines=center,
xlabel style={at={(current axis.right of origin)},xshift=1.5ex,anchor=center},
ylabel style={at={(current axis.above origin)},yshift=1.5ex, anchor=center}}}
\pgfkeys{/pgfplots/grafikh parastash/.style={red!80!black,line width=.4mm,samples=200}}

\newlist{alist}{enumerate}{3}
\setlist[alist]{itemsep=0mm,label=\alph*.}

\newcommand{\kerkissans}[1]{{\fontfamily{maksf}\selectfont \textbf{#1}}}
\renewcommand{\textdexiakeraia}{}
\newcommand{\dint}{\displaystyle\int}

%----------- ΑΣΚΗΣΗ ------------------
\newcommand\themanum{A}
\newcounter{thema}[section]
\renewcommand{\thethema}{\arabic{thema}}   
\newcommand{\Thema}{\refstepcounter{thema}{\textbf{\textcolor{black}{\faPenSquare\ \  \large \kerkissans{\textbf{\thethemaο\hspace{2mm}Θέμα}}}}}\hspace{0mm}}{}

\newlist{erwthma}{enumerate}{3}
\setlist[erwthma]{label=\bf{\kerkissans{\large{\textcolor{black}{\themanum\arabic*}}}},itemsep=0mm,leftmargin=0.8cm}

\newenvironment{thema}[1]
{\begin{tcolorbox}[title=\Thema\kerkissans{\textbf{\large{  #1}}},breakable,
enhanced standard,titlerule=0pt,toprule=0pt, rightrule=0pt, bottomrule=0pt,
colback=white,left=0mm,top=1mm,bottom=0mm,
boxrule=0pt,
colframe=white,leftrule=0mm,sharp corners,coltitle=black]
\renewcommand{\themanum}{#1}}
{\end{tcolorbox}}
%-----------------------------------------


\begin{document}

\title{\vspace{-2cm}\kerkissans{\textbf{Επαναληπτικά Θέματα Γ Λυκείου - \underline{10 Πλήρη Θέματα Α}}}\vspace{-0.5cm}}
\date{}
\maketitle

% =================== ΔΙΑΓΩΝΙΣΜΑ 1 ===================
\begin{thema}{A}
\begin{erwthma}
    \item Έστω ένα πολυώνυμο $P(x) = a_\nu x^\nu + a_{\nu-1} x^{\nu-1} + \ldots + a_1 x + a_0$ και $ x_0 \in \mathbb{R} $. Να αποδείξετε ότι ισχύει: 
    \[ \lim_{x\to x_0} P(x) = P(x_0) \]
    \item Πότε μια συνάρτηση $f:A\to\mathbb{R}$ ονομάζεται συνάρτηση $1-1$;
    \item Πώς ορίζεται η αντίστροφη συνάρτηση $f^{-1}$ μιας $1-1$ συνάρτησης $f$;
    \item Να χαρακτηρίσετε τις παρακάτω προτάσεις ως Σωστές (Σ) ή Λανθασμένες (Λ).
    \begin{alist}
        \item Αν δυο συναρτήσεις $f, g$ έχουν το ίδιο πεδίο ορισμού και κοινό σύνολο τιμών, τότε είναι ίσες.
        \item Το $\lim_{x\to 0} \frac{\cos x - 1}{x} = 0$.
        \item Αν μια συνάρτηση είναι 1-1, αυτό συνεπάγεται ότι είναι γνησίως μονότονη.
        \item Το πεδίο ορισμού της $f \circ g$ είναι πάντα υποσύνολο του πεδίου ορισμού της $g$.
        \item Για τη σταθερή συνάρτηση $f(x)=c, c\in\mathbb{R}$ ισχύει $f'(x)=c$.
    \end{alist}
    \item Έστω οι συναρτήσεις $f, g$. Το πεδίο ορισμού της σύνθεσης $f \circ g$ ορίζεται ως:
    \begin{tasks}[label=\let\textdexiakeraia\relax\alph*.](2)
        \task $\big\{x \in D_f \mid g(x) \in D_g\big\}$
        \task $\big\{x \in D_g \mid f(x) \in D_f\big\}$
        \task $\big\{x \in D_g \mid g(x) \in D_f\big\}$
        \task $\big\{x \in D_f \mid f(x) \in D_g\big\}$
    \end{tasks}
\end{erwthma}
\end{thema}
\vspace{1cm}

% =================== ΔΙΑΓΩΝΙΣΜΑ 2 ===================
\begin{thema}{A}
\begin{erwthma}
    \item Να διατυπώσετε το Θεώρημα του Bolzano και να δώσετε τη γεωμετρική του ερμηνεία με ένα σχήμα.
    \item Πότε μια συνάρτηση $f$ λέγεται παραγωγίσιμη σε ένα ανοικτό διάστημα $(a, \beta)$ και πότε σε ένα κλειστό $[a, \beta]$;
    \item Πότε μια συνάρτηση $f$ ονομάζεται συνεχής;
    \item Να χαρακτηρίσετε τις παρακάτω προτάσεις ως Σωστές (Σ) ή Λανθασμένες (Λ).
    \begin{alist}
        \item Αν $f$ συνεχής στο $[\alpha, \beta]$ και $f(\alpha)f(\beta)>0$, τότε η εξίσωση $f(x)=0$ αποκλείεται να έχει ρίζα στο $(\alpha, \beta)$.
        \item Το $\lim_{x\to \infty} \frac{\sin x}{x} = 1$.
        \item Αν η $f$ δεν είναι συνεχής στο $x_0$, τότε σίγουρα δεν είναι παραγωγίσιμη στο $x_0$.
        \item Αν $\lim_{x\to x_0} f(x) = +\infty$, τότε υποχρεωτικά $\lim_{x\to x_0} \frac{1}{f(x)} = 0$.
        \item Αν μια συνάρτηση δεν είναι 1-1, δεν αντιστρέφεται.
    \end{alist}
    \item Αν $\lim_{x\to x_0} f(x) = 0$ και $\lim_{x\to x_0} g(x) = +\infty$, τότε το $\lim_{x\to x_0} [f(x)g(x)]$ είναι:
    \begin{tasks}[label=\let\textdexiakeraia\relax\alph*.](4)
        \task Πάντα $0$
        \task Πάντα $+\infty$
        \task Πάντα $-\infty$
        \task Απροσδιόριστη μορφή
    \end{tasks}
\end{erwthma}
\end{thema}
\vspace{1cm}

% =================== ΔΙΑΓΩΝΙΣΜΑ 3 ===================
\setcounter{page}{1} % Just to keep numbering clean if pages break
\begin{thema}{A}
\begin{erwthma}
    \item Να αποδείξετε ότι αν μια συνάρτηση $f$ είναι παραγωγίσιμη σε ένα σημείο $x_0$, τότε είναι και συνεχής στο σημείο αυτό.
    \item Τι ονομάζουμε ολικό μέγιστο και τι ολικό ελάχιστο μιας συνάρτησης $f$;
    \item Τι ονομάζεται γραφική παράσταση μιας συνάρτησης $f$;
    \item Να χαρακτηρίσετε τις παρακάτω προτάσεις ως Σωστές (Σ) ή Λανθασμένες (Λ).
    \begin{alist}
        \item Αν η $f$ είναι συνεχής στο $x_0$, τότε είναι υποχρεωτικά και παραγωγίσιμη στο $x_0$.
        \item Η συνάρτηση $f(x) = |x|$ δεν είναι παραγωγίσιμη στο $x_0=0$.
        \item Αν $\int_\alpha^\beta f(x) dx \ge 0$, τότε αναγκαστικά $f(x) \ge 0$ για κάθε $x \in [\alpha, \beta]$.
        \item Av $\lim_{x\to x_0} f(x) = -\infty$, τότε για τιμές κοντά στο $x_0$ ισχύει $f(x)<0$.
        \item Το θεώρημα Fermat ισχύει και για ακρότατα που βρίσκονται στα άκρα γραφικής παράστασης (κλειστό διάστημα).
    \end{alist}
    \item Η παράγωγος της ολοκληρωτικής συνάρτησης $F(x) = \dint_\alpha^x f(t) dt$ (με $f$ συνεχή) είναι:
    \begin{tasks}[label=\let\textdexiakeraia\relax\alph*.](4)
        \task $F'(x) = 0$
        \task $F'(x) = f(t)$
        \task $F'(x) = f(x)$
        \task $F'(x) = f(x) - f(\alpha)$
    \end{tasks}
\end{erwthma}
\end{thema}
\vspace{1cm}

% =================== ΔΙΑΓΩΝΙΣΜΑ 4 ===================
\begin{thema}{A}
\begin{erwthma}
    \item Έστω δύο παραγωγίσιμες συναρτήσεις $f, g$. Να αποδείξετε, με χρήση του ορισμού, ότι η συνάρτηση $f+g$ είναι παραγωγίσιμη και ισχύει $(f+g)'(x) = f'(x) + g'(x)$.
    \item Πότε η ευθεία $x=x_0$ ονομάζεται κατακόρυφη ασύμπτωτη της γραφικής παράστασης μιας συνάρτησης $f$;
    \item Πότε η ευθεία $y=y_0$ ονομάζεται οριζόντια ασύμπτωτη της γραφικής παράστασης μιας συνάρτησης $f$;
    \item Να χαρακτηρίσετε τις παρακάτω προτάσεις ως Σωστές (Σ) ή Λανθασμένες (Λ).
    \begin{alist}
        \item Αν δυο συναρτήσεις $f, g$ τέμνονται πάνω στην $y=x$, τότε αυτές είναι απαραίτητα αντίστροφες.
        \item Αν $\lim f(x)=+\infty$ και $\lim g(x)=-\infty$, τότε πάντα $\lim[f(x)+g(x)]=0$.
        \item Στο θεώρημα Rolle, αν ελαττώσουμε τον περιορισμό της συνέχειας στα κλειστά άκρα, το συμπέρασμα ενδέχεται να μην ισχύει.
        \item Το ολοκλήρωμα $\int_a^\beta f(x) dx$ εξαρτάται από τον συμβολισμό της μεταβλητής ολοκλήρωσης (π.χ. $x$ αντί για $t$).
        \item Αν $f''(x) > 0$ για κάθε $x \in \Delta$, τότε η $f$ στρέφει τα κοίλα προς τα άνω.
    \end{alist}
    \item Ο γεωμετρικός ορισμός της $f'(x_0)$ ταυτίζεται με:
    \begin{tasks}[label=\let\textdexiakeraia\relax\alph*.](2)
        \task Τον άξονα συμμετρίας της συνάρτησης.
        \task Το εμβαδόν υπό την καμπύλη.
        \task Το συντελεστή διεύθυνσης $\lambda$ της εφαπτομένης της $C_f$ στο $x_0$.
        \task Την ελάχιστη τιμή της $f$.
    \end{tasks}
\end{erwthma}
\end{thema}
\vspace{1cm}

% =================== ΔΙΑΓΩΝΙΣΜΑ 5 ===================
\begin{thema}{A}
\begin{erwthma}
    \item Να διατυπώσετε το Θεώρημα του Rolle και να το ερμηνεύσετε γεωμετρικά με τη βοήθεια κατάλληλου σχήματος.
    \item Τι ονομάζουμε ρυθμό μεταβολής ενός ποσού $y=f(x)$ ως προς ένα ποσό $x$ σε ένα σημείο $x_0$;
    \item Πότε λέμε ότι μια συνάρτηση $f$ στρέφει τα κοίλα προς τα άνω ή είναι κυρτή στο $\Delta$;
    \item Να χαρακτηρίσετε τις παρακάτω προτάσεις ως Σωστές (Σ) ή Λανθασμένες (Λ).
    \begin{alist}
        \item Ο ρυθμός μεταβολής της ταχύτητας $v(t)$ ισούται με την επιτάχυνση $\alpha(t)$.
        \item Αν η γραφική παράσταση της $f$ βρίσκεται εξ ολοκλήρου "κάτω" από κάθε εφαπτομένη της, τότε είναι κυρτή.
        \item Η γραφική παράσταση της $f(x) = \sin x$ παρουσιάζει άπειρα σημεία καμπής.
        \item Κάθε σταθερή συνάρτηση είναι παραγωγίσιμη με $f'(x)=0$.
        \item Το πεδίο ορισμού του $\ln(x^2 - 1)$ είναι το $(1, +\infty)$.
    \end{alist}
    \item Αν $f''(x) > 0$ για κάθε $x \in \Delta$, η γραφική παράσταση της $f$ αναφορικά με την εφαπτομένη της:
    \begin{tasks}[label=\let\textdexiakeraia\relax\alph*.](2)
        \task Βρίσκεται χαμηλότερα, με εξαίρεση το σημείο επαφής.
        \task Βρίσκεται ψηλότερα, με εξαίρεση το σημείο επαφής.
        \task Διαπερνάται από αυτή σε κάθε σημείο.
        \task Συμπίπτει με αυτή.
    \end{tasks}
\end{erwthma}
\end{thema}
\vspace{1cm}

% =================== ΔΙΑΓΩΝΙΣΜΑ 6 ===================
\begin{thema}{A}
\begin{erwthma}
    \item Να αποδείξετε ότι η παράγωγος της ταυτοτικής συνάρτησης $f(x) = x$ είναι ίση με $1$, δηλαδή $(x)' = 1$.
    \item Πότε το σημείο $A(x_0, f(x_0))$ λέγεται σημείο καμπής της γραφικής παράστασης της $f$;
    \item Πώς ορίζονται οι πράξεις της πρόσθεσης $f+g$ και του πολλαπλασιασμού $f \cdot g$ δύο συναρτήσεων; Ποιο το πεδίο ορισμού τους;
    \item Να χαρακτηρίσετε τις παρακάτω προτάσεις ως Σωστές (Σ) ή Λανθασμένες (Λ).
    \begin{alist}
        \item Αν η συνάρτηση $f$ έχει συντηρητικό πρόσημο $f(x)>0$ και είναι συνεχής, τότε $\int_a^\beta f(x)dx=0$ (με $a<\beta$).
        \item Η συνάρτηση $\tan x$ δεν παραγωγίζεται στα σημεία $x = \kappa \pi + \frac{\pi}{2}, \kappa \in \mathbb{Z}$.
        \item Το θεώρημα ενδιάμεσων τιμών εφαρμόζεται πάντα και σε ανοιχτά διαστήματα $(a,\beta)$.
        \item Σημείο καμπής υπάρχει μόνο στα σημεία όπου $f''(x_0)=0$.
        \item Μια συνάρτηση 1-1 δεν μπορεί να είναι άρτια.
    \end{alist}
    \item Το $A(x_0, f(x_0))$ λέγεται σημείο καμπής αν:
    \begin{tasks}[label=\let\textdexiakeraia\relax\alph*.](2)
        \task H $f$ παρουσιάζει τοπικό μέγιστο και τοπικό ελάχιστο.
        \task H $C_f$ έχει εφαπτομένη και αλλάζει κυρτότητα εκατέρωθεν του $x_0$.
        \task $f'(x_0) = 0$ αυστηρά.
        \task $f(x_0) = 0$.
    \end{tasks}
\end{erwthma}
\end{thema}
\vspace{1cm}

% =================== ΔΙΑΓΩΝΙΣΜΑ 7 ===================
\begin{thema}{A}
\begin{erwthma}
    \item Έστω μια συνάρτηση $f$ συνεχής σε ένα διάστημα $\Delta$. Αν $f'(x)>0$ σε κάθε εσωτερικό σημείο του $\Delta$, να αποδείξετε ότι η συνάρτηση $f$ είναι γνησίως αύξουσα σε όλο το $\Delta$.
    \item Πότε η ευθεία $y=\lambda x+\beta$ ονομάζεται πλάγια ασύμπτωτη της γραφικής παράστασης μιας συνάρτησης $f$ στο $+\infty$;
    \item Τι ονομάζουμε αρχική συνάρτηση ή παράγουσα μιας συνάρτησης $f$ σε ένα διάστημα $\Delta$;
    \item Να χαρακτηρίσετε τις παρακάτω προτάσεις ως Σωστές (Σ) ή Λανθασμένες (Λ).
    \begin{alist}
        \item Αν $f'(x)>0$ για κάθε $x \in \mathbb{R}^*$, τότε η $f$ είναι υποχρεωτικά γνησίως αύξουσα στο $\mathbb{R}^*$.
        \item Αν σε μια συνεχής συνάρτηση ισχύει $\int_\alpha^\beta f(x) dx = 0$, τότε υποχρεωτικά η $f(x)=0$ παντού.
        \item Κάθε κατακόρυφη ασύμπτωτη δημιουργείται από ρίζα του παρονομαστή που δεν απλοποιείται.
        \item Είναι εφικτό μια συνάρτηση να έχει και οριζόντια και πλάγια ασύμπτωτη στο $+\infty$.
        \item Το θεώρημα μέγιστης και ελάχιστης τιμής εξασφαλίζει ότι η συνάρτηση παίρνει τις τιμές αυτές σίγουρα στο εσωτερικό του $[a, \beta]$.
    \end{alist}
    \item Έστω $f(x) \ge 0$. Το εμβαδόν του χωρίου μεταξύ $C_f$, άξονα $x'x$ και ευθειών $x=\alpha, x=\beta$ (με $\alpha<\beta$) δίνεται από:
    \begin{tasks}[label=\let\textdexiakeraia\relax\alph*.](4)
        \task $E = \dint_\beta^\alpha f(x) dx$
        \task $E = - \dint_\alpha^\beta f(x) dx$
        \task $E = \dint_\alpha^\beta f(x) dx$
        \task $E = |f(\alpha) - f(\beta)|$
    \end{tasks}
\end{erwthma}
\end{thema}
\vspace{1cm}

% =================== ΔΙΑΓΩΝΙΣΜΑ 8 ===================
\begin{thema}{A}
\begin{erwthma}
    \item Να διατυπώσετε το Θεώρημα Μέσης Τιμής (Θ.Μ.Τ.) του διαφορικού λογισμού. Τι εκφράζει το θεώρημα αυτό γεωμετρικά;
    \item Τι ονομάζουμε πραγματική συνάρτηση;
    \item Τι ονομάζουμε τοπικό ελάχιστο και τι ολικό ελάχιστο μιας συνάρτησης;
    \item Να χαρακτηρίσετε τις παρακάτω προτάσεις ως Σωστές (Σ) ή Λανθασμένες (Λ).
    \begin{alist}
        \item Η γραφική παράσταση της συνάρτησης -f είναι συμμετρική της f ως προς τον άξονα x'x.
        \item Αν $f''(x) = 0$ στο $x_0$, τότε το $x_0$ είναι υποχρεωτικά σημείο καμπής.
        \item Η αντίστροφη μιας γνησίως φθίνουσας συνάρτησης είναι επίσης γνησίως φθίνουσα.
        \item Το εμβαδόν μεταξύ της $f(x) = \cos x$ και του άξονα x'x στο $[0, 2\pi]$ ισούται με 0.
        \item Το Θ.Μ.Τ αφορά εφαπτομένη η οποία είναι παράλληλη στη χορδή $A(\alpha, f(\alpha))$ και $B(\beta, f(\beta))$.
    \end{alist}
    \item Αν $\lim_{x\to x_0} f(x) = -\infty$, τότε για τιμές του $x$ "πολύ κοντά" στο $x_0$, ισχύει:
    \begin{tasks}[label=\let\textdexiakeraia\relax\alph*.](4)
        \task $f(x) > 0$
        \task $f(x) < 0$
        \task $f(x) = 0$
        \task $\lim f(x) \in \mathbb{R}$
    \end{tasks}
\end{erwthma}
\end{thema}
\vspace{1cm}

% =================== ΔΙΑΓΩΝΙΣΜΑ 9 ===================
\begin{thema}{A}
\begin{erwthma}
    \item Έστω $f(x) = x^\nu$ όπου $\nu$ φυσικός αριθμός. Με τη βοήθεια του ορισμού της παραγώγου να αποδείξετε ότι η $(x^\nu)' = \nu x^{\nu - 1}$. 
    \item Πότε δύο συναρτήσεις $f, g$ ονομάζονται ίσες;
    \item Πότε μια συνάρτηση λέγεται γνησίως φθίνουσα στο διάστημα $\Delta$;
    \item Να χαρακτηρίσετε τις παρακάτω προτάσεις ως Σωστές (Σ) ή Λανθασμένες (Λ).
    \begin{alist}
        \item Οι συναρτήσεις $f(x) = \ln(x^2)$ και $g(x) = 2\ln x$ είναι ίσες στο $\mathbb{R}^*$.
        \item Αν η $f'$ μηδενίζεται σε ένα σημείο $x_0$, τότε παρουσιάζει εκεί πάντα ακρότατο.
        \item $\dint f(x)g'(x) dx = f(x)g(x) - \dint f'(x)g(x) dx$.
        \item $\lim_{x \to -\infty} e^x = 0$.
        \item Για $x > 0$ ισχύει ότι η συνάρτηση $\ln x$ έχει κατακόρυφη ασύμπτωτη τον άξονα y'y.
    \end{alist}
    \item Ο τύπος ολοκλήρωσης κατά παράγοντες είναι:
    \begin{tasks}[label=\let\textdexiakeraia\relax\alph*.](2)
        \task $\dint_\alpha^\beta f(x)g'(x) dx = f(\beta)g(\beta) - \dint_\alpha^\beta f'(x)g(x) dx$
        \task $\dint_\alpha^\beta f(x)g'(x) dx = [f(x)g(x)]_\alpha^\beta - \dint_\alpha^\beta f'(x)g(x) dx$
        \task $\dint_\alpha^\beta f'(x)g'(x) dx = f(x)g(x)$
        \task $\dint f g' dx = f g$
    \end{tasks}
\end{erwthma}
\end{thema}
\vspace{1cm}

% =================== ΔΙΑΓΩΝΙΣΜΑ 10 ===================
\begin{thema}{A}
\begin{erwthma}
    \item Να αποδείξετε ότι οι γραφικές παραστάσεις $ C_f $ και $ C_{f^{-1}} $ των συναρτήσεων $ f $ και $ f^{-1} $ είναι συμμετρικές ως προς την ευθεία $ y=x $ (διχοτόμο $x\hat{O}y$).
    \item Τι ονομάζεται σύνολο τιμών μιας πραγματικής συνάρτησης $f$;
    \item Τι λέγεται τοπικό μέγιστο της $f$;
    \item Να χαρακτηρίσετε τις παρακάτω προτάσεις ως Σωστές (Σ) ή Λανθασμένες (Λ).
    \begin{alist}
        \item Δεν μπορούμε να εφαρμόσουμε τον κανόνα L'Hospital σε μορφές $\frac{0}{0}$ ή $\frac{\infty}{\infty}$ αν δεν υπάρχει το όριο των παραγώγων.
        \item Ένα τοπικό ελάχιστο μπορεί να είναι μεγαλύτερο από ένα τοπικό μέγιστο της ίδιας συνάρτησης.
        \item Κάθε συνεχής συνάρτηση σε διάστημα $[a,\beta]$ είναι ενσωματώσιμη, και έχει πάντα αρχική συνάρτηση $F(x) = \dint_\alpha^x f(t) dt$.
        \item Η γραφική παράσταση της $|f(x)|$ προκύπτει αν αναδιπλώσουμε τα αρνητικά $y$ πάνω από τον άξονα $x'$.
        \item Το θεώρημα ενδιάμεσων τιμών (Θ.Ε.Τ) ταυτίζεται ουσιαστικά με το θεώρημα της μέγιστης / ελάχιστης τιμής.
    \end{alist}
    \item Αν $x_0$ είναι εσωτερικό σημείο και τοπικό ακρότατο της $f$, και η $f$ παραγωγίζεται στο $x_0$, τότε:
    \begin{tasks}[label=\let\textdexiakeraia\relax\alph*.](4)
        \task $f'(x_0) > 0$
        \task $f'(x_0) < 0$
        \task $f'(x_0) = 0$
        \task Δεν εξάγεται συμπέρασμα
    \end{tasks}
\end{erwthma}
\end{thema}

\end{document}
